\chapter[SVD \& Positive Definite Matrices]{Singular Value Decomposition (SVD), Low-Rank Approximations \& Positive Definite Matrices}
\label{chap:week06}

% ----------------------------------------------------------------------
% PART 1: EXECUTIVE ROADMAP & LEARNING OBJECTIVES
% ----------------------------------------------------------------------
\section{Executive Roadmap \& Learning Objectives}

While eigendecomposition is strictly confined to square matrices and can fail entirely if a matrix is defective, the \textbf{Singular Value Decomposition (SVD)} stands as the undisputed crown jewel of linear algebra: it exists for \emph{every} matrix $A \in \mathbb{R}^{m \times n}$, whether square or rectangular, full-rank or rank-deficient, well-conditioned or completely singular. 

SVD unlocks the deepest geometric truth of linear transformations: \textbf{every linear map transforms a unit sphere into a hyper-ellipsoid}. SVD dissects this geometric transformation into three elementary, decoupled physical actions: a rigid rotation in the domain, a coordinate scaling along perpendicular axes by non-negative \textbf{singular values}, and a final rigid rotation in the codomain.

In data science and modern machine learning, SVD is the foundation of dimensionality reduction (Principal Component Analysis), matrix compression, latent semantic analysis, recommender systems (collaborative filtering), and regularized pseudo-inversion. This chapter develops the complete mathematical theory of SVD, the Eckart-Young-Mirsky low-rank approximation theorem, and the characterization of positive definite matrices.

\subsection{Prerequisites \& Scaffolding}
Mastery of this chapter builds directly upon orthogonal matrices, fundamental subspaces, and the Spectral Theorem as mapped in Table~\ref{tab:w06_prereqs}.

\begin{table}[htbp]
\centering
\small
\begin{tabularx}{\textwidth}{lXX}
\toprule
\textbf{Foundational Area} & \textbf{Core Concepts Required} & \textbf{Role in Machine Learning Foundations} \\
\midrule
Spectral Theorem & Orthogonal diagonalization $S = Q \Lambda Q^T$ of symmetric matrices & Foundation for decomposing Gram matrices $A^T A$ and $A A^T$ \\
Fundamental Subspaces & $\Col(A), \Row(A), \Null(A), \Null(A^T)$, dimensions and orthogonality & Explicitly mapped to the singular vectors $U$ and $V$ \\
Matrix Norms & Frobenius norm $\|A\|_F$, spectral operator norm $\|A\|_2$ & Quantifying error in low-rank data compression \\
\bottomrule
\end{tabularx}
\caption{Prerequisite Scaffolding for Chapter 6.}
\label{tab:w06_prereqs}
\end{table}

\subsection{Core Competencies}
By the conclusion of this chapter, students will be able to:
\begin{itemize}
    \item Derive and compute both the full SVD ($A = U \Sigma V^T$) and compact SVD ($A = U_r \Sigma_r V_r^T$) for arbitrary rectangular matrices.
    \item Relate singular values and singular vectors directly to the eigendecompositions of the Gram matrices $A^T A$ and $A A^T$.
    \item State and apply the four equivalent characterizations of Positive Definite and Semidefinite matrices (eigenvalues, quadratic forms, Sylvester's minors, and Cholesky decomposition).
    \item State and apply the Eckart-Young-Mirsky Theorem to compute the provably optimal rank-$k$ matrix approximation under both Spectral and Frobenius norms.
    \item Construct the Moore-Penrose pseudoinverse $A^+$ via SVD to find the unique minimum-norm solution to arbitrary linear systems.
\end{itemize}

\subsection{Notation \& Mathematical Conventions}
\begin{itemize}
    \item $A \in \mathbb{R}^{m \times n}$: An arbitrary real matrix of rank $r \le \min(m, n)$.
    \item $\sigma_1 \ge \sigma_2 \ge \dots \ge \sigma_r > 0$: The strictly positive singular values of $A$, ordered decreasingly.
    \item $U = [\mathbf{u}_1, \dots, \mathbf{u}_m] \in \mathbb{R}^{m \times m}$: Orthogonal matrix of left singular vectors (eigenvectors of $A A^T$).
    \item $V = [\mathbf{v}_1, \dots, \mathbf{v}_n] \in \mathbb{R}^{n \times n}$: Orthogonal matrix of right singular vectors (eigenvectors of $A^T A$).
    \item $\Sigma \in \mathbb{R}^{m \times n}$: Rectangular pseudo-diagonal matrix with singular values $\Sigma_{ii} = \sigma_i$.
    \item $A_k \triangleq \sum_{i=1}^k \sigma_i \mathbf{u}_i \mathbf{v}_i^T$: The optimal truncated rank-$k$ approximation of $A$.
    \item $A^+ \triangleq V \Sigma^+ U^T = \sum_{i=1}^r \frac{1}{\sigma_i} \mathbf{v}_i \mathbf{u}_i^T$: The Moore-Penrose pseudoinverse.
    \item $M \succ 0$ ($M \succeq 0$): Denotes a symmetric positive definite (positive semidefinite) matrix.
\end{itemize}

% ----------------------------------------------------------------------
% PART 2: THE GRAND PARADOX & MOTIVATION
% ----------------------------------------------------------------------
\section{The Grand Paradox: The Sphere-to-Ellipsoid Transformation}

\subsection{The Universal Geometric Action}
Consider an arbitrary $2 \times 2$ matrix with messy, non-symmetric entries:
\begin{equation}
A = \begin{bmatrix} 3 & 0 \\ 4 & 5 \end{bmatrix}
\end{equation}
If you feed a circle of unit vectors $\mathbf{x} \in \mathbb{R}^2$ ($\|\mathbf{x}\|_2 = 1$) into this matrix, what does the output set of vectors $\{A\mathbf{x} \mid \|\mathbf{x}\|_2 = 1\}$ look like? Does it warp into a peanut shape? Does it fold into a polygon?

\begin{quote}
\textit{\textbf{The Geometric Discovery:}} \textbf{Every} linear transformation—without exception—maps a unit sphere into a hyper-ellipsoid. 
\end{quote}

Why? Because if $\mathbf{y} = A\mathbf{x}$ and $A$ is invertible, then $\mathbf{x} = A^{-1} \mathbf{y}$. The unit sphere condition $\|\mathbf{x}\|_2^2 = 1$ becomes:
\begin{equation}
\mathbf{x}^T \mathbf{x} = 1 \implies (A^{-1}\mathbf{y})^T (A^{-1}\mathbf{y}) = 1 \implies \mathbf{y}^T (A A^T)^{-1} \mathbf{y} = 1
\end{equation}
The matrix $(A A^T)^{-1}$ is strictly symmetric positive definite! By the Spectral Theorem, any quadratic equation $\mathbf{y}^T M \mathbf{y} = 1$ with $M \succ 0$ defines an \textbf{ellipsoid}. 

The principal axes of this output ellipsoid are the left singular vectors $\mathbf{u}_i$, the half-lengths of the axes are the singular values $\sigma_i$, and the original orthogonal directions on the unit sphere that get stretched into these axes are the right singular vectors $\mathbf{v}_i$!

\begin{ponderbox}[Singular Values vs. Eigenvalues]
Consider the matrix $A = \begin{bmatrix} 0 & 3 \\ -3 & 0 \end{bmatrix}$.
\begin{enumerate}
    \item Compute the eigenvalues of $A$. Are they real or complex?
    \item Compute the Gram matrix $A^T A$ and find its eigenvalues.
    \item What are the singular values of $A$? Can a singular value ever be negative or complex?
    \item Why does SVD succeed in extracting real scale factors when eigendecomposition yields purely imaginary numbers?
\end{enumerate}
\end{ponderbox}

% ----------------------------------------------------------------------
% PART 3: MATHEMATICAL THEORY OF SVD
% ----------------------------------------------------------------------
\section{Mathematical Theory of SVD}

\begin{theorem}[Singular Value Decomposition (SVD)]
Let $A \in \mathbb{R}^{m \times n}$ have rank $r \le \min(m, n)$. There exist orthogonal matrices $U \in \mathbb{R}^{m \times m}$ ($U^T U = I_m$) and $V \in \mathbb{R}^{n \times n}$ ($V^T V = I_n$), and a unique diagonal matrix $\Sigma \in \mathbb{R}^{m \times n}$ such that:
\begin{equation}
A = U \Sigma V^T = \sum_{i=1}^r \sigma_i \mathbf{u}_i \mathbf{v}_i^T
\label{eq:svd_outer}
\end{equation}
where the singular values are strictly positive and ordered:
\begin{equation}
\sigma_1 \ge \sigma_2 \ge \dots \ge \sigma_r > 0
\end{equation}
\end{theorem}

\begin{proof}
Consider the Gram matrix $A^T A \in \mathbb{R}^{n \times n}$.
\begin{enumerate}
    \item \textbf{Symmetry and Semidefiniteness:}
    $(A^T A)^T = A^T (A^T)^T = A^T A$, so $A^T A$ is symmetric.
    For any $\mathbf{x} \in \mathbb{R}^n$, $\mathbf{x}^T (A^T A)\mathbf{x} = (A\mathbf{x})^T (A\mathbf{x}) = \|A\mathbf{x}\|_2^2 \ge 0$. Thus $A^T A$ is positive semidefinite.
    \item \textbf{Orthonormal Basis $V$:}
    By the Spectral Theorem for real symmetric matrices, $A^T A$ possesses an orthonormal basis of eigenvectors $\mathbf{v}_1, \dots, \mathbf{v}_n \in \mathbb{R}^n$ with non-negative real eigenvalues:
    \begin{equation}
    A^T A \mathbf{v}_i = \lambda_i \mathbf{v}_i, \quad \lambda_1 \ge \lambda_2 \ge \dots \ge \lambda_n \ge 0
    \end{equation}
    Because $\rank(A^T A) = \rank(A) = r$, exactly $r$ eigenvalues are strictly positive: $\lambda_i > 0$ for $i \le r$, and $\lambda_i = 0$ for $i > r$.
    \item \textbf{Definition of Singular Values:}
    Define $\sigma_i \triangleq \sqrt{\lambda_i} > 0$ for $i \in \{1, \dots, r\}$.
    \item \textbf{Construction of Orthonormal Vectors $U$:}
    For each $i \in \{1, \dots, r\}$, define:
    \begin{equation}
    \mathbf{u}_i \triangleq \frac{1}{\sigma_i} A \mathbf{v}_i \in \mathbb{R}^m
    \label{eq:def_ui}
    \end{equation}
    Testing the inner product between $\mathbf{u}_i$ and $\mathbf{u}_j$:
    \begin{align}
    \mathbf{u}_i^T \mathbf{u}_j &= \left( \frac{1}{\sigma_i} A \mathbf{v}_i \right)^T \left( \frac{1}{\sigma_j} A \mathbf{v}_j \right) = \frac{1}{\sigma_i \sigma_j} \mathbf{v}_i^T (A^T A \mathbf{v}_j) \\
    &= \frac{1}{\sigma_i \sigma_j} \mathbf{v}_i^T (\lambda_j \mathbf{v}_j) = \frac{\lambda_j}{\sigma_i \sigma_j} \mathbf{v}_i^T \mathbf{v}_j = \frac{\sigma_j^2}{\sigma_i \sigma_j} \delta_{ij} = \delta_{ij}
    \end{align}
    Thus $\{\mathbf{u}_1, \dots, \mathbf{u}_r\}$ forms an orthonormal set in $\mathbb{R}^m$. Extend this set via Gram-Schmidt to an orthonormal basis $\{\mathbf{u}_1, \dots, \mathbf{u}_m\}$ for all of $\mathbb{R}^m$.
    \item \textbf{Matrix Assembly:}
    From Equation~\eqref{eq:def_ui}, $A\mathbf{v}_i = \sigma_i \mathbf{u}_i$ for $i \le r$. For $i > r$, $\lambda_i = 0 \implies \|A\mathbf{v}_i\|_2^2 = \mathbf{v}_i^T A^T A \mathbf{v}_i = 0 \implies A\mathbf{v}_i = \mathbf{0}$.
    Stacking into matrices:
    \begin{equation}
    A V = U \Sigma \implies A = U \Sigma V^T
    \end{equation}
\end{enumerate}
\end{proof}

\begin{figure}[htbp]
\centering
\begin{tikzpicture}[
    scale=0.9,
    font=\sffamily\footnotesize
]

% 1. Input Space: Unit Circle
\draw[thick, fill=primarylight!30] (-5.0, 0) circle (1.5cm);
\draw[->, >=stealth, color=gray] (-6.8, 0) -- (-3.2, 0) node[right] {$x_1$};
\draw[->, >=stealth, color=gray] (-5.0, -1.8) -- (-5.0, 1.8) node[above] {$x_2$};
\draw[->, >=stealth, line width=1.5pt, color=primarydark] (-5.0, 0) -- (-3.8, 0.9) node[above right] {$\mathbf{v}_1$};
\draw[->, >=stealth, line width=1.5pt, color=primarydark] (-5.0, 0) -- (-5.9, 1.2) node[above left] {$\mathbf{v}_2$};
\node[below] at (-5.0, -2.0) {\textbf{1. Domain: } Basis $V$};

% First arrow
\draw[->, >=stealth, line width=1.5pt, color=primarydark] (-3.0, 0) -- (-1.5, 0) node[midway, above] {$V^T$};

% 2. Intermediate Space: Scaled Ellipse aligned with axes
\draw[thick, fill=tealbg!30] (0.5, 0) ellipse (2.0cm and 0.9cm);
\draw[->, >=stealth, color=gray] (-1.8, 0) -- (2.8, 0);
\draw[->, >=stealth, color=gray] (0.5, -1.5) -- (0.5, 1.5);
\draw[->, >=stealth, line width=1.5pt, color=tealborder] (0.5, 0) -- (2.5, 0) node[below] {$\sigma_1 \mathbf{e}_1$};
\draw[->, >=stealth, line width=1.5pt, color=tealborder] (0.5, 0) -- (0.5, 0.9) node[left] {$\sigma_2 \mathbf{e}_2$};
\node[below] at (0.5, -2.0) {\textbf{2. Stretched: } $\Sigma$};

% Second arrow
\draw[->, >=stealth, line width=1.5pt, color=primarydark] (3.0, 0) -- (4.5, 0) node[midway, above] {$U$};

% 3. Output Space: Rotated Ellipse
\begin{scope}[shift={(6.5, 0)}, rotate=40]
    \draw[thick, fill=coralbg!40] (0, 0) ellipse (2.0cm and 0.9cm);
    \draw[->, >=stealth, line width=1.5pt, color=coralborder] (0, 0) -- (2.0, 0) node[above right] {$\sigma_1 \mathbf{u}_1$};
    \draw[->, >=stealth, line width=1.5pt, color=coralborder] (0, 0) -- (0, 0.9) node[above left] {$\sigma_2 \mathbf{u}_2$};
\end{scope}
\node[below] at (6.5, -2.0) {\textbf{3. Codomain: } Basis $U$};

\end{tikzpicture}
\caption{The Geometric Action of SVD: Any matrix transformation factorizes into a domain rotation $V^T$, axis stretching $\Sigma$, and codomain rotation $U$.}
\label{fig:svd_geometry}
\end{figure}

% ----------------------------------------------------------------------
% PART 4: POSITIVE DEFINITE MATRICES
% ----------------------------------------------------------------------
\section{Positive Definite \& Semidefinite Matrices}

\begin{definitionbox}{Positive Definite and Semidefinite Matrices}
A real symmetric matrix $M \in \mathbb{R}^{n \times n}$ ($M = M^T$) is:
\begin{itemize}
    \item \textbf{Positive Definite} ($M \succ 0$) if for all $\mathbf{x} \in \mathbb{R}^n \setminus \{\mathbf{0}\}$: $\mathbf{x}^T M \mathbf{x} > 0$.
    \item \textbf{Positive Semidefinite} ($M \succeq 0$) if for all $\mathbf{x} \in \mathbb{R}^n$: $\mathbf{x}^T M \mathbf{x} \ge 0$.
\end{itemize}
\end{definitionbox}

\begin{theorem}[Equivalent Tests for Positive Definiteness]
For any real symmetric matrix $M \in \mathbb{R}^{n \times n}$, the following four conditions are mathematically equivalent:
\begin{enumerate}
    \item \textbf{Eigenvalue Test:} All eigenvalues are strictly positive: $\lambda_i > 0$ for all $i \in \{1, \dots, n\}$.
    \item \textbf{Energy Test:} The quadratic form $\mathbf{x}^T M \mathbf{x} > 0$ for all non-zero $\mathbf{x} \in \mathbb{R}^n$.
    \item \textbf{Sylvester's Leading Minors Test:} All $n$ leading principal minors are strictly positive:
    \begin{equation}
    \det(M_{1:k, 1:k}) > 0 \quad \text{for } k = 1, 2, \dots, n
    \end{equation}
    \item \textbf{Cholesky Decomposition Test:} There exists a unique lower-triangular matrix $L$ with strictly positive diagonal entries ($L_{ii} > 0$) such that:
    \begin{equation}
    M = L L^T
    \label{eq:cholesky}
    \end{equation}
\end{enumerate}
\end{theorem}

\begin{proof}[Proof Sketch]
$(1 \iff 2)$: By the Spectral Theorem, $M = Q \Lambda Q^T$. Let $\mathbf{y} = Q^T \mathbf{x}$. Then $\mathbf{x}^T M \mathbf{x} = \mathbf{y}^T \Lambda \mathbf{y} = \sum_{i=1}^n \lambda_i y_i^2$. Because $Q$ is invertible, $\mathbf{x} \ne \mathbf{0} \iff \mathbf{y} \ne \mathbf{0}$. The sum $\sum \lambda_i y_i^2$ is strictly positive for all $\mathbf{y} \ne \mathbf{0}$ if and only if every $\lambda_i > 0$.

$(1 \iff 4)$: If $M = L L^T$, then for $\mathbf{x} \ne \mathbf{0}$: $\mathbf{x}^T M \mathbf{x} = \mathbf{x}^T L L^T \mathbf{x} = \|L^T \mathbf{x}\|_2^2$. Since $L$ has positive diagonals, $\det(L) = \prod L_{ii} > 0$, so $L$ is invertible. Hence $L^T \mathbf{x} \ne \mathbf{0}$ whenever $\mathbf{x} \ne \mathbf{0}$, ensuring $\|L^T \mathbf{x}\|_2^2 > 0$.
\end{proof}

% ----------------------------------------------------------------------
% PART 5: LOW-RANK APPROXIMATION & ECKART-YOUNG
% ----------------------------------------------------------------------
\section{Low-Rank Matrix Approximation \& Eckart-Young Theorem}

In massive machine learning datasets (user-item rating matrices in Netflix, document-word matrices in NLP, or high-resolution images), storing a full matrix $A \in \mathbb{R}^{m \times n}$ requires $\mathcal{O}(mn)$ numbers. Can we compress $A$ into a matrix $B$ of low rank $k \ll \min(m, n)$ while discarding as little information as possible?

\begin{theorem}[Eckart-Young-Mirsky Theorem]
Let $A = \sum_{i=1}^r \sigma_i \mathbf{u}_i \mathbf{v}_i^T$ be the SVD of $A \in \mathbb{R}^{m \times n}$ with $\sigma_1 \ge \sigma_2 \ge \dots \ge \sigma_r > 0$. For any target rank $k < r$, the truncated SVD:
\begin{equation}
A_k \triangleq \sum_{i=1}^k \sigma_i \mathbf{u}_i \mathbf{v}_i^T
\label{eq:truncated_svd_formula}
\end{equation}
is the \textbf{globally optimal} rank-$k$ approximation of $A$ under both the Spectral (operator) norm and Frobenius norm:
\begin{align}
\min_{\rank(B) \le k} \|A - B\|_2 &= \|A - A_k\|_2 = \sigma_{k+1} \label{eq:ey_spectral_error} \\
\min_{\rank(B) \le k} \|A - B\|_F &= \|A - A_k\|_F = \sqrt{\sum_{i=k+1}^r \sigma_i^2} \label{eq:ey_frobenius_error}
\end{align}
\end{theorem}

\begin{ponderbox}[The Low-Rank Compression Miracle]
Consider a $1000 \times 1000$ image ($1,000,000$ pixels).
\begin{enumerate}
    \item How many numbers are required to store the full image $A$?
    \item If you approximate the image with a rank-$k = 20$ truncated SVD $A_{20} = U_{20} \Sigma_{20} V_{20}^T$, how many numbers do you need to store (columns of $U$, singular values, columns of $V$)?
    \item What is the compression ratio?
    \item Why is this compression provably superior to any other rank-20 matrix that could ever be created?
\end{enumerate}
\end{ponderbox}

% ----------------------------------------------------------------------
% PART 6: IN-DEPTH WORKED EXAMPLES
% ----------------------------------------------------------------------
\section{In-Depth Worked Numerical Examples}

\begin{examplebox}{Manual SVD of a Rectangular $3 \times 2$ Matrix}
Let matrix $A \in \mathbb{R}^{3 \times 2}$ be defined as:
\begin{equation}
A = \begin{bmatrix} 1 & 1 \\ 0 & 1 \\ -1 & 1 \end{bmatrix}
\end{equation}
Compute the full SVD $A = U \Sigma V^T$ step-by-step by hand.

\textbf{Step 1: Compute Gram Matrix $A^T A$ and Eigendecomposition.}
\begin{equation}
A^T A = \begin{bmatrix} 1 & 0 & -1 \\ 1 & 1 & 1 \end{bmatrix} \begin{bmatrix} 1 & 1 \\ 0 & 1 \\ -1 & 1 \end{bmatrix} = \begin{bmatrix} 1+0+1 & 1+0-1 \\ 1+0-1 & 1+1+1 \end{bmatrix} = \begin{bmatrix} 2 & 0 \\ 0 & 3 \end{bmatrix}
\end{equation}
Notice $A^T A$ is already diagonal!
The eigenvalues are $\lambda_1 = 3$ and $\lambda_2 = 2$ (ordered decreasingly).
The singular values are:
\begin{equation}
\sigma_1 = \sqrt{3}, \quad \sigma_2 = \sqrt{2}
\end{equation}

\textbf{Step 2: Determine Orthonormal Right Singular Vectors $V$.}
Because $A^T A = \begin{bmatrix} 2 & 0 \\ 0 & 3 \end{bmatrix}$:
\begin{itemize}
    \item For $\lambda_1 = 3$: $(A^T A - 3I)\mathbf{v}_1 = \begin{bmatrix} -1 & 0 \\ 0 & 0 \end{bmatrix}\mathbf{v}_1 = \mathbf{0} \implies \mathbf{v}_1 = \begin{bmatrix} 0 \\ 1 \end{bmatrix}$.
    \item For $\lambda_2 = 2$: $(A^T A - 2I)\mathbf{v}_2 = \begin{bmatrix} 0 & 0 \\ 0 & 1 \end{bmatrix}\mathbf{v}_2 = \mathbf{0} \implies \mathbf{v}_2 = \begin{bmatrix} 1 \\ 0 \end{bmatrix}$.
\end{itemize}
Thus:
\begin{equation}
V = \begin{bmatrix} \mathbf{v}_1 & \mathbf{v}_2 \end{bmatrix} = \begin{bmatrix} 0 & 1 \\ 1 & 0 \end{bmatrix}
\end{equation}

\textbf{Step 3: Compute Left Singular Vectors $\mathbf{u}_1, \mathbf{u}_2 \in \mathbb{R}^3$.}
Using $\mathbf{u}_i = \frac{1}{\sigma_i} A \mathbf{v}_i$:
\begin{align}
\mathbf{u}_1 &= \frac{1}{\sqrt{3}} \begin{bmatrix} 1 & 1 \\ 0 & 1 \\ -1 & 1 \end{bmatrix} \begin{bmatrix} 0 \\ 1 \end{bmatrix} = \frac{1}{\sqrt{3}} \begin{bmatrix} 1 \\ 1 \\ 1 \end{bmatrix} \\
\mathbf{u}_2 &= \frac{1}{\sqrt{2}} \begin{bmatrix} 1 & 1 \\ 0 & 1 \\ -1 & 1 \end{bmatrix} \begin{bmatrix} 1 \\ 0 \end{bmatrix} = \frac{1}{\sqrt{2}} \begin{bmatrix} 1 \\ 0 \\ -1 \end{bmatrix}
\end{align}
Testing orthogonality: $\mathbf{u}_1^T \mathbf{u}_2 = \frac{1}{\sqrt{6}}(1(1) + 1(0) + 1(-1)) = 0 \quad \checkmark$.

\textbf{Step 4: Find the Third Basis Vector $\mathbf{u}_3 \in \Null(A^T)$.}
To construct full $U \in \mathbb{R}^{3 \times 3}$, we need $\mathbf{u}_3$ orthogonal to both $\mathbf{u}_1$ and $\mathbf{u}_2$. Taking cross product or solving $A^T \mathbf{u}_3 = \mathbf{0}$:
\begin{equation}
\begin{bmatrix} 1 & 0 & -1 \\ 1 & 1 & 1 \end{bmatrix} \begin{bmatrix} z_1 \\ z_2 \\ z_3 \end{bmatrix} = \begin{bmatrix} 0 \\ 0 \end{bmatrix} \implies z_1 = z_3, \quad z_2 = -2z_3
\end{equation}
Choosing $z_3 = 1 \implies [1, -2, 1]^T$. Normalizing:
\begin{equation}
\mathbf{u}_3 = \frac{1}{\sqrt{6}} \begin{bmatrix} 1 \\ -2 \\ 1 \end{bmatrix}
\end{equation}

\textbf{Step 5: Assemble Full SVD and Reconstruct.}
\begin{equation}
U = \begin{bmatrix} 1/\sqrt{3} & 1/\sqrt{2} & 1/\sqrt{6} \\ 1/\sqrt{3} & 0 & -2/\sqrt{6} \\ 1/\sqrt{3} & -1/\sqrt{2} & 1/\sqrt{6} \end{bmatrix}, \quad \Sigma = \begin{bmatrix} \sqrt{3} & 0 \\ 0 & \sqrt{2} \\ 0 & 0 \end{bmatrix}, \quad V^T = \begin{bmatrix} 0 & 1 \\ 1 & 0 \end{bmatrix}
\end{equation}
Reconstructing $U \Sigma V^T$:
\begin{align}
U \Sigma &= \begin{bmatrix} 1 & 1 \\ 1 & 0 \\ 1 & -1 \end{bmatrix} \\
U \Sigma V^T &= \begin{bmatrix} 1 & 1 \\ 1 & 0 \\ 1 & -1 \end{bmatrix} \begin{bmatrix} 0 & 1 \\ 1 & 0 \end{bmatrix} = \begin{bmatrix} 1 & 1 \\ 0 & 1 \\ -1 & 1 \end{bmatrix} = A \quad \checkmark
\end{align}
\end{examplebox}

\begin{examplebox}{Positive Definiteness Diagnostic: Sylvester vs. Cholesky}
Let matrix $M = \begin{bmatrix} 2 & -1 & 0 \\ -1 & 2 & -1 \\ 0 & -1 & 2 \end{bmatrix}$.
\begin{enumerate}
    \item Compute all leading principal minors (Sylvester's criterion).
    \item Construct the lower-triangular Cholesky factor $L$ ($M = L L^T$).
    \item Expand the quadratic form $\mathbf{x}^T M \mathbf{x}$ as a sum of squares.
\end{enumerate}

\textbf{Step 1: Compute Leading Principal Minors.}
\begin{align}
\Delta_1 &= \det[2] = 2 > 0 \\
\Delta_2 &= \det\begin{bmatrix} 2 & -1 \\ -1 & 2 \end{bmatrix} = 4 - 1 = 3 > 0 \\
\Delta_3 &= \det(M) = 2(4 - 1) - (-1)(-2 - 0) + 0 = 6 - 2 = 4 > 0
\end{align}
Since $\Delta_1 = 2, \Delta_2 = 3, \Delta_3 = 4 > 0$, Sylvester's criterion guarantees $M \succ 0$.

\textbf{Step 2: Cholesky Factorization $M = L L^T$.}
Let $L = \begin{bmatrix} l_{11} & 0 & 0 \\ l_{21} & l_{22} & 0 \\ l_{31} & l_{32} & l_{33} \end{bmatrix}$. Matching entries of $L L^T = M$:
\begin{align}
l_{11}^2 = 2 &\implies l_{11} = \sqrt{2} \\
l_{21} l_{11} = -1 &\implies l_{21} = -\frac{1}{\sqrt{2}} \\
l_{21}^2 + l_{22}^2 = 2 &\implies \frac{1}{2} + l_{22}^2 = 2 \implies l_{22} = \sqrt{\frac{3}{2}} \\
l_{31} l_{11} = 0 &\implies l_{31} = 0 \\
l_{31} l_{21} + l_{32} l_{22} = -1 &\implies 0 + l_{32} \sqrt{\frac{3}{2}} = -1 \implies l_{32} = -\sqrt{\frac{2}{3}} \\
l_{31}^2 + l_{32}^2 + l_{33}^2 = 2 &\implies 0 + \frac{2}{3} + l_{33}^2 = 2 \implies l_{33} = \sqrt{\frac{4}{3}} = \frac{2}{\sqrt{3}}
\end{align}
The Cholesky factor is:
\begin{equation}
L = \begin{bmatrix} \sqrt{2} & 0 & 0 \\ -1/\sqrt{2} & \sqrt{3/2} & 0 \\ 0 & -\sqrt{2/3} & 2/\sqrt{3} \end{bmatrix}
\end{equation}

\textbf{Step 3: Quadratic Form as Sum of Squares.}
\begin{align}
\mathbf{x}^T M \mathbf{x} &= 2x_1^2 + 2x_2^2 + 2x_3^2 - 2x_1 x_2 - 2x_2 x_3 \\
&= x_1^2 + (x_1^2 - 2x_1 x_2 + x_2^2) + (x_2^2 - 2x_2 x_3 + x_3^2) + x_3^2 \\
&= x_1^2 + (x_1 - x_2)^2 + (x_2 - x_3)^2 + x_3^2
\end{align}
Because this is a sum of four squares, it is strictly non-negative. It equals zero if and only if $x_1 = 0$, $x_1 = x_2$, $x_2 = x_3$, and $x_3 = 0 \implies \mathbf{x} = \mathbf{0}$. 

Hence $\mathbf{x}^T M \mathbf{x} > 0$ for all $\mathbf{x} \ne \mathbf{0}$, providing independent visual proof that $M \succ 0$!
\end{examplebox}

% ----------------------------------------------------------------------
% PART 7: EXAM TRAPS & DIAGNOSTIC MISCONCEPTIONS
% ----------------------------------------------------------------------
\section{Exam Traps \& Diagnostic Misconceptions}

\begin{trapbox}{Gotcha 1: Conflating Singular Values with Eigenvalues}
A student claims that because a matrix has eigenvalues $\lambda_1 = 5, \lambda_2 = -3$, its singular values must be $5$ and $-3$.
\textbf{The Trap:} Singular values are the square roots of the eigenvalues of $A^T A$: $\sigma_i = \sqrt{\lambda_i(A^T A)}$. Because $A^T A \succeq 0$, singular values are strictly \textbf{non-negative real numbers} ($\sigma_i \ge 0$). They never take negative or complex values. For a real symmetric matrix with negative eigenvalues, the singular values are the absolute values $|\lambda_i|$.
\end{trapbox}

\begin{trapbox}{Gotcha 2: Believing All-Positive Matrix Entries Imply Positive Definiteness}
A programmer checks a covariance candidate matrix $C = \begin{bmatrix} 1 & 3 \\ 3 & 1 \end{bmatrix}$ and assumes it is positive definite because all entries are positive ($C_{ij} > 0$).
\textbf{The Trap:} Positive matrix entries have \textbf{no relation} to positive definiteness. Evaluating the determinant: $\det(C) = 1(1) - 3(3) = -8 < 0$. The eigenvalues are $\lambda_1 = 4$ and $\lambda_2 = -2$. For $\mathbf{x} = [1, -1]^T$, the quadratic form is $\mathbf{x}^T C \mathbf{x} = 1 - 3 - 3 + 1 = -4 < 0$. The matrix is indefinite!
\end{trapbox}

\begin{trapbox}{Gotcha 3: Assuming $U$ and $V$ are Interchangeable}
An analyst switches $U$ and $V$, writing $A = V \Sigma U^T$.
\textbf{The Trap:} $U \in \mathbb{R}^{m \times m}$ contains the eigenvectors of $A A^T$ (spanning $\mathbb{R}^m$), while $V \in \mathbb{R}^{n \times n}$ contains the eigenvectors of $A^T A$ (spanning $\mathbb{R}^n$). If $A$ is rectangular ($m \ne n$), swapping $U$ and $V$ yields a matrix of dimension $n \times m$, which is the transpose $A^T$, not $A$.
\end{trapbox}

\begin{trapbox}{Gotcha 4: Measuring Variance Retained by Linear Sum of Singular Values}
In truncated SVD, an engineer claims that rank-$k$ approximation retains $\frac{\sum_{i=1}^k \sigma_i}{\sum_{i=1}^r \sigma_i}$ of the total dataset variance.
\textbf{The Trap:} Total variance (Frobenius energy) corresponds to the sum of \textbf{squared} singular values: $\|A\|_F^2 = \sum_{i=1}^r \sigma_i^2$. The true proportion of retained energy is $\frac{\sum_{i=1}^k \sigma_i^2}{\sum_{i=1}^r \sigma_i^2}$. Using linear sums severely overestimates the information preserved in weak singular components.
\end{trapbox}

% ----------------------------------------------------------------------
% PART 8: COMPARATIVE SYNTHESIS TABLE
% ----------------------------------------------------------------------
\section{Comparative Synthesis Table}

Table~\ref{tab:eigen_vs_svd} provides a comprehensive comparison between Eigendecomposition and Singular Value Decomposition.

\begin{table}[htbp]
\centering
\small
\begin{tabularx}{\textwidth}{lXX}
\toprule
\textbf{Property} & \textbf{Eigendecomposition ($A = P D P^{-1}$)} & \textbf{Singular Value Decomposition ($A = U \Sigma V^T$)} \\
\midrule
Matrix Applicability & Square matrices only ($n \times n$) & Any rectangular matrix ($m \times n$) \\
\addlinespace
Existence Guarantee & Only if matrix is non-defective & Always exists for every matrix \\
\addlinespace
Basis Vectors & Columns of $P$ (generally non-orthogonal) & Columns of $U, V$ (strictly orthonormal) \\
\addlinespace
Diagonal Elements & Eigenvalues $\lambda_i \in \mathbb{C}$ (can be negative/complex) & Singular values $\sigma_i \in \mathbb{R}_{\ge 0}$ (sorted decreasingly) \\
\addlinespace
Geometric Action & Stretches along invariant eigenspaces & Rotates domain ($V^T$), scales axes ($\Sigma$), rotates codomain ($U$) \\
\addlinespace
Low-Rank Truncation & Generally not optimal for low-rank approximation & Provably optimal under both $L_2$ and Frobenius norms \\
\bottomrule
\end{tabularx}
\caption{Comparative Synthesis: Eigendecomposition vs. Singular Value Decomposition.}
\label{tab:eigen_vs_svd}
\end{table}

% ----------------------------------------------------------------------
% PART 9: PYTHON / SCIENTIFIC COMPUTING LAB
% ----------------------------------------------------------------------
\section{Python / Scientific Computing Lab}

The production script below implements truncated SVD for matrix compression, verifies the Eckart-Young-Mirsky theorem, and validates positive definiteness via Cholesky decomposition.

\begin{lstlisting}[language=Python, caption={SVD Low-Rank Approximation and Positive Definiteness in Python.}]
import numpy as np

def compute_svd_approximation(A: np.ndarray, k: int):
    """Computes truncated rank-k SVD approximation: A_k = sum_{i=1}^k sigma_i u_i v_i^T."""
    U, s, Vt = np.linalg.svd(A, full_matrices=False)
    # Truncate to top k components
    U_k = U[:, :k]
    s_k = s[:k]
    Vt_k = Vt[:k, :]
    A_k = U_k @ np.diag(s_k) @ Vt_k

    # Compute theoretical vs empirical Frobenius error
    emp_frob_err = np.linalg.norm(A - A_k, 'fro')
    theory_frob_err = np.sqrt(np.sum(s[k:] ** 2))
    return A_k, emp_frob_err, theory_frob_err

def check_positive_definite(M: np.ndarray) -> bool:
    """Checks positive definiteness using Cholesky factorization."""
    # Symmetrize to guard against numerical float errors
    M_sym = (M + M.T) / 2.0
    try:
        _ = np.linalg.cholesky(M_sym)
        return True
    except np.linalg.LinAlgError:
        return False

def main():
    rng = np.random.RandomState(42)
    m, n = 50, 20
    # Generate low-rank matrix with noise
    rank_true = 5
    A_clean = rng.randn(m, rank_true) @ rng.randn(rank_true, n)
    noise = 0.05 * rng.randn(m, n)
    A = A_clean + noise

    print("Verifying Eckart-Young-Mirsky Theorem:")
    for k in [1, 3, 5, 10]:
        _, emp_err, theory_err = compute_svd_approximation(A, k)
        print(f"  Rank-{k:2d}: Empirical ||A - A_k||_F = {emp_err:.4f}, Theoretical = {theory_err:.4f}")

    # Positive Definite Checks
    M_pos = A.T @ A
    M_indef = np.array([[1.0, 3.0], [3.0, 1.0]])

    print(f"\nA^T A is Positive Definite: {check_positive_definite(M_pos)}")
    print(f"Matrix [[1, 3], [3, 1]] is Positive Definite: {check_positive_definite(M_indef)}")

if __name__ == "__main__":
    main()
\end{lstlisting}

% ----------------------------------------------------------------------
% PART 10: COMPREHENSIVE PRACTICE PROBLEMS
% ----------------------------------------------------------------------
\section{Comprehensive Practice Problems}

\subsection*{Part A: Analytical \& Mathematical Problems}

\begin{enumerate}[label=\textbf{P\arabic*.}]
    \item \textbf{Explicit Full SVD Computation by Hand:}
    Let $A = \begin{bmatrix} 3 & 0 \\ 4 & 5 \end{bmatrix}$.
    \begin{enumerate}[label=(\alph*)]
        \item Compute the Gram matrix $A^T A$ and find its eigenvalues $\lambda_1, \lambda_2$.
        \item Determine the singular values $\sigma_1, \sigma_2$.
        \item Find the orthonormal right singular vectors $\mathbf{v}_1, \mathbf{v}_2$.
        \item Compute left singular vectors $\mathbf{u}_i = \frac{1}{\sigma_i} A \mathbf{v}_i$ and verify $U \Sigma V^T = A$.
    \end{enumerate}

    \item \textbf{Sylvester's Criterion vs. Cholesky Factorization:}
    Let $M = \begin{bmatrix} 2 & -1 & 0 \\ -1 & 2 & -1 \\ 0 & -1 & 2 \end{bmatrix}$.
    \begin{enumerate}[label=(\alph*)]
        \item Compute all three leading principal minors: $\det(M_{1:1})$, $\det(M_{1:2, 1:2})$, and $\det(M)$.
        \item Use Sylvester's Criterion to determine if $M$ is strictly positive definite.
        \item Construct the lower-triangular Cholesky factor $L$ such that $M = L L^T$.
    \end{enumerate}

    \item \textbf{Frobenius Norm Representation in Terms of Singular Values:}
    Prove that for any matrix $A \in \mathbb{R}^{m \times n}$ with singular values $\sigma_1, \dots, \sigma_r$:
    \begin{equation}
    \|A\|_F^2 = \trace(A^T A) = \sum_{i=1}^r \sigma_i^2
    \end{equation}

    \item \textbf{Condition Number and Sensitivity in Linear Inversion:}
    The condition number of a matrix $A$ is defined as $\kappa(A) = \frac{\sigma_{\max}}{\sigma_{\min}}$.
    Suppose $A\mathbf{x} = \mathbf{b}$ and $A(\mathbf{x} + \delta \mathbf{x}) = \mathbf{b} + \delta \mathbf{b}$. Prove that:
    \begin{equation}
    \frac{\|\delta \mathbf{x}\|_2}{\|\mathbf{x}\|_2} \le \kappa(A) \frac{\|\delta \mathbf{b}\|_2}{\|\mathbf{b}\|_2}
    \end{equation}

    \item \textbf{Pseudoinverse of Rank-Deficient Matrix:}
    Let $A = \begin{bmatrix} 1 & 1 \\ 2 & 2 \end{bmatrix}$.
    \begin{enumerate}[label=(\alph*)]
        \item Determine $\rank(A)$ and compute its compact SVD $A = \sigma_1 \mathbf{u}_1 \mathbf{v}_1^T$.
        \item Compute the Moore-Penrose pseudoinverse $A^+ = \frac{1}{\sigma_1} \mathbf{v}_1 \mathbf{u}_1^T$.
        \item Verify that $A A^+ A = A$ and $A^+ A A^+ = A^+$.
    \end{enumerate}
\end{enumerate}

\subsection*{Part B: Diagnostic Multiple-Choice Questions}

\begin{enumerate}[label=\textbf{MCQ\arabic*.}]
    \item Let $A \in \mathbb{R}^{20 \times 10}$ have singular values $\sigma_1 = 12, \sigma_2 = 8, \sigma_3 = 5, \sigma_4 = 2$, with all other singular values equal to zero. What is the minimum Frobenius norm error $\min_{\rank(B) \le 2} \|A - B\|_F$?
    \begin{enumerate}[label=(\alph*)]
        \item $\sqrt{5^2 + 2^2} = \sqrt{29}$
        \item $5$
        \item $5 + 2 = 7$
        \item $\sqrt{12^2 + 8^2} = \sqrt{208}$
    \end{enumerate}

    \item Which of the following conditions is necessary and sufficient for a real symmetric matrix $A$ to be positive definite?
    \begin{enumerate}[label=(\alph*)]
        \item All entries $A_{ij} > 0$.
        \item The trace $\trace(A) > 0$ and determinant $\det(A) > 0$.
        \item All eigenvalues of $A$ are strictly positive.
        \item The diagonal entries are strictly greater than off-diagonal entries.
    \end{enumerate}
\end{enumerate}
